\begin{textsample}{POS Dim 4 – ce – Score 15.00 – ce/ce\_em\_8.txt}  \label{ex:f4_pos_131}
2.
Sujeitos do Ensino Médio O Documento Curricular Referencial do Ceará ( DCRC ) é resultado dos debates e da reorganização curricular promovidos pela Base Nacional Comum Curricular ( BNCC ).
Portanto, alguns aspectos devem ser considerados ao curso da construção de todo e qualquer currículo que se origine neste documento, sejam em escolas públicas ou particulares do \textbf{estado} do Ceará.
Um desses aspectos corresponde às/aos cidadãs/ãos a quem se destina a finalidade do currículo.
Posto isso, o DCRC deve conduzir uma organização curricular que vise ao respeito à \textbf{identidade} dos diferentes tipos de sujeitos presentes no Ensino Médio, com foco na formação integral e humana, desenvolvendo -os a nível pessoal, social e laboral.
É necessário, dessa forma, reconhecer as diversidades de sujeitos em suas várias características e necessidades, sejam elas de \textbf{território}, de cultura, de faixa etária, de etnia, de religião, de condição socioeconômica, de gênero etc.
Por isso, ao longo deste capítulo, serão abordadas as diferentes concepções dos sujeitos potencialmente presentes no Ensino Médio cearense, através de uma visão contemporânea, buscando a promoção, nos espaços escolares, da identificação, da aceitação e da integração dos estudantes, tornando a escola um ambiente seguro para se expor e se propor enquanto “ sujeitos inconclusos”², proporcionando a possibilidade de “ ser mais ”, na perspectiva de Freire ( 2019 ). ² Paulo Freire afirma que os seres humanos são dotados de um constante ampliar de si, ou seja, ao curso da vida, a cada momento, estamos aptos a apreender ou a construir, a inventar, a modificar pensamentos, ações de modo a estarmos em permanente estágio de ampliação, podendo “ ser mais ”, melhor dizendo, ampliando o seu ser a um plano mais.
Logo, cabe às/aos atrizes/atores envolvidas/os na construção do currículo escolar a reflexão, a produção e a criação de caminhos para um Ensino Médio que contemple os objetivos de uma educação qualitativa, igualitária e equitativa, principalmente voltada a uma aprendizagem significativa para seus sujeitos. 2.1 As/os Jovens e as Juventudes Conforme o Estatuto da Juventude ( Lei nº 12.852/13 ), são consideradas/os jovens pessoas com idade entre 15 e 29 anos.
Dado o exposto, é importante lembrar que, neste recorte etário, ainda se encontram jovens com idade entre 15 e 18 anos também assistidos pelo Estatuto da Criança e do Adolescente ( Lei nº 8.069/90 ).
Portanto, compreende -se que a maior parte dos sujeitos atendidos pelo Ensino Médio encontra -se nesse perfil etário e que qualquer estudo e planejamento para construção curricular desta etapa do ensino necessita de uma melhor compreensão acerca do público juvenil.
Entende -se a juventude como um período de transição biopsicossocial situado entre a infância e a fase adulta, em que, ao longo de seu desenvolvimento, os indivíduos se encontram em aspiração à plena existência, como também no desejo por uma vida simultaneamente autônoma e comunitária.
Neste momento, principalmente no recorte etário correspondente à adolescência ( entre 15 e 18 anos – ECA/90 ), o indivíduo se encontra em uma catarse emocional, recheada de esperanças e desesperanças, de transgressões, de dúvidas, de angústias e de sonhos.
Por isso, é característica a esse período a necessidade de adaptação e acomodação psicológica e social.
Ao trabalhar com a juventude, deve -se compreender as diversas facetas e os elementos que atuam diretamente sobre o desenvolvimento biopsicossocial destes indivíduos, o que provoca perfis bem diversos e democráticos, exigindo a aceitação da pluralidade que o termo nos confere : juventudes.
Nessa perspectiva, cabe à escola oferecer oportunidades únicas para o desenvolvimento de competências que possibilitem o exercício de uma cidadania plena, a propriedade no domínio digital, assim como a preparação para o moderno mundo do trabalho, abordando habilidades, atitudes e valores que contribuam para a resolução de demandas complexas da vida cotidiana.
Portanto, é imprescindível que se reconheça o jovem como pessoa, como sujeito social e como cidadã/ão, criando espaços para que ela/ele possa compreender a si mesma/o ( e à/ao outra/o ), descobrir -se e expressar o seu próprio potencial. 2.2 As/os Adultas/os e as/os Idosas/os As escolas atendem a demandas de alunas/os de faixas etárias distintas, sendo, portanto, de fundamental importância que voltemos nossa atenção, também, para o público adulto e idoso, uma vez que muitos, por inúmeros fatores, não tiveram a oportunidade de realizar seus estudos básicos na idade adequada.
Em concordância com o Estatuto da Criança e do Adolescente ( ECA ), assim como com o Estatuto do Idoso ( EI ), são considerados adultos todos os indivíduos com idade entre dezoito e cinquenta e nove anos, enquanto que todos aqueles com idade igual ou superior aos sessenta anos são considerados idosos.
Conforme Rutter ( 1996 ), “ é um equívoco acreditar que as aprendizagens sejam exclusivas às crianças e aos adolescentes. ” Sendo assim, todos os adultos e todos as/os idosas/os estão aptos a aprender, a mudar e a evoluir.
É preponderante, portanto, que busquemos compreender as especificidades que \textbf{carregam} as/os alunas/os dessas faixas etárias, percebendo que muitos deles estão distantes da prática escolar, desabituados aos ritos escolares, envoltos de uma história de vida própria e com objetivos específicos para sua formação.
Precisamos, então, pensar em estratégias e práticas educacionais mais fidedignas ao contexto delas/deles.
Diante dessas constatações, é de responsabilidade da comunidade escolar o acolhimento dessas/es estudantes por meio de uma prática educacional humana.
Compreendemos que os estímulos nessas faixas etárias precisam estar conectados à vida real, respeitando o ritmo com que se aprende, aproximando os aspectos estudados às necessidades cotidianas delas/es, garantindo sua permanência na escola e, consequentemente, a conclusão do Ensino Médio.
Em síntese, as escolas cearenses têm o papel de acolher os diversos públicos de forma plena, considerando as necessidades e os objetivos que emanam de suas/seus estudantes.
Significa, assim, que devem proporcionar uma visão pedagógica que se adeque ao contexto da educação de adultas/os e de idosas/os, garantindo, além da socialização e da (re)inclusão dessas pessoas no universo escolar e na sociedade, o direito à aprendizagem plena das competências e das habilidades determinadas pela BNCC, assim como uma adequação metodológica conforme as faixas etárias distintas. 2.3 As Mulheres Mesmo diante do histórico de conquistas das \textbf{mulheres}, principalmente a partir do fortalecimento do Movimento Feminista na década de 1970 e com a própria criação do conceito de gênero, pensado não como algo natural, mas construído histórica e socialmente, cabe pensar que ainda são imensos os desafios atuais impostos ao mundo feminino, no qual a \textbf{mulher} tem que assumir inúmeros papéis sociais, muitas vezes de forma desigual.
As transformações sociais, ocorridas desde o século passado, provocaram mudanças no modo de agir de \textbf{homens} e de \textbf{mulheres}, o que alterou a forma como ambos se comportam na sociedade, quer no âmbito familiar, quer no profissional.
Essas mudanças se relacionam diretamente com o conceito de gênero, a partir do qual as diferenças biológicas não definem os papéis sociais, sendo estes construídos culturalmente.
Essa compreensão afirma que há uma diversidade de modos de ser \textbf{mulher}, fato observado também no contexto escolar.
Contudo, essas mudanças na forma das \textbf{mulheres} estarem no mundo não anularam todos os discursos que buscam dar a elas funções específicas, como secretária, zeladora, merendeira, dentre outras ; muitas delas com menor importância social e mais laboriosas que as delegadas aos \textbf{homens}.
Apesar das conquistas dos movimentos sociais e feministas, ainda prevalece uma desigualdade que \textbf{hierarquiza} os gêneros, dando aos \textbf{homens}, muitas vezes, lugares privilegiados em relação às \textbf{mulheres}.
Essas construções sociais se estendem ao ambiente escolar, configurando -se como um desafio ao acesso, à permanência, à conclusão dos estudos e à inserção de muitas \textbf{mulheres} no mercado de trabalho.
Consideramos importante, também, refletir e discutir sobre a necessidade de combater, a partir da educação e dentro da escola, o conceito de patriarcado, pois \textbf{mulheres}, desde meninas, passam por uma complexidade de situações através de uma relação de gênero que envolve “ um sistema de dominação que se faz presente nas diferentes instituições sociais, desde a família ao \textbf{Estado}, apresentando -se em todos os espaços da sociedade ” ( ALMEIDA, 2017, p. 56 ). 2.4 As Negras e os Negros Em decorrência de processos históricos que fomentaram, sobretudo, a escravidão dos povos africanos, articulados com discursos que visavam justificar a violência, a crueldade e a injustiça contra eles, várias nações foram constituídas tendo o \textbf{racismo} como elemento estrutural.
A sociedade brasileira também (re)produz \textbf{racismo} e, mesmo diante de muitas lutas dos povos negros, percebemos que ele está presente no espaço escolar.
É nesse cenário que indicamos que a mudança passa, sobretudo, pela criação de políticas públicas específicas, bem como pela educação.
É, então, a partir das lutas desses movimentos que leis foram criadas para fomentar o respeito à diversidade \textbf{racial} e garantir a esses sujeitos o acesso e a permanência nos processos educacionais.
Referimo -nos aqui às leis que dão origem à política de Educação para as Relações Étnico-Raciais ( ERER ), Lei nº 10.639, de 09 de janeiro de 2003, e Lei nº 11.645, de 10 de março de 2008, que tornam obrigatório o ensino de história e cultura africana, afro-brasileira e \textbf{indígena} em toda a extensão do currículo da Educação Básica.
Entretanto, para além de temáticas no currículo, essas leis tratam mais profundamente do acesso, da permanência e do sucesso das/os alunas/os negras e negros na escola, e isso passa pela valorização de suas existências e de sua história.
Assim, o \textbf{estado} do Ceará defende que a efetivação da cidadania passa, também, pelo reconhecimento das matrizes culturais e de conhecimento que constituem as diversidades do Brasil, sendo a escola uma das primeiras instituições que tem a função de formar cidadãs/ãos na perspectiva do respeito, da valorização da diferença \textbf{racial}, étnica e de gênero.
Dessa forma, as culturas e os conhecimentos negros possibilitam às/aos negras/os e às/aos não negras/os a criação de uma nova história e de novas \textbf{identidades}, fortalecendo e valorizando a positividade, a beleza, a radicalidade e a presença africana em nossa formação cultural.
Desse modo, as/os alunas/os precisam ter acesso ao conhecimento produzido pelos povos negros ( africanos, brasileiros e cearenses ) a fim de que tenham condições reais de compreender sua história para enfrentar o \textbf{racismo}. 2.5 Povos \textbf{Indígenas} De acordo com o Estatuto do \textbf{Índio}, é considerado \textbf{indígena} “ todo indivíduo de origem e ascendência pré-colombiana que se identifica e é identificado como pertencente a um grupo étnico cujas características culturais o distinguem da sociedade nacional ” ( BRASIL, 1973 ).
Associada a essa percepção, o Decreto nº 10.088, de 05 de novembro de 2019, construído a partir da Convenção 169 da OIT ( Organização Internacional do Trabalho ), declara que povos \textbf{indígenas} são aqueles que descendem “ de \textbf{populações} que habitavam o país ou uma região geográfica pertencente ao país na época da conquista ou da colonização ou do estabelecimento das atuais \textbf{fronteiras} estatais e que, seja qual for sua situação jurídica, conservam todas as suas próprias instituições sociais, econômicas, culturais e políticas, ou parte delas ”.
Em outras palavras, a FUNAI ( Fundação Nacional do \textbf{Índio} ) sinaliza para uma construção da \textbf{identidade} étnica \textbf{indígena} com base na autodeclaração do indivíduo e o reconhecimento dessa \textbf{identidade} por parte do grupo de origem.
Nesse sentido, diversas/os alunas/os \textbf{indígenas} passaram a pertencer a comunidades integradas, assim chamadas aquelas que, conforme o estatuto, foram incorporadas à sociedade não \textbf{indígena}.
Tais comunidades ainda lutam pelo reconhecimento de seus direitos civis e humanos, conservando os usos, os costumes e as tradições característicos da sua cultura.
Essas/es estudantes devem ser respeitadas/os e valorizadas/os como cidadãs/ãos conscientes de sua \textbf{identidade} cultural.
São jovens que, considerando sua cultura, devem ser preparadas/os, não apenas para enfrentarem os desafios e vivenciarem as demandas do século XXI, mas, principalmente, para conquistar seus lugares na universidade, na política, nos negócios, na ciência etc., a fim de fortalecer sua comunidade, seu \textbf{território}, seu povo, suas \textbf{identidades} e suas culturas.
Diante disso, as ações executadas no Ceará para o desenvolvimento da educação desses povos devem se alicerçar nas Diretrizes Curriculares Nacionais para Educação Escolar \textbf{Indígena} e no Plano de Ação do \textbf{Território} Etnoeducacional Potirõ, em regime de colaboração entre os sistemas de ensino e demais instituições envolvidas com a modalidade, respeitando sua \textbf{territorialidade}.
A essas/es estudantes é assegurado o direito de aprender sobre e na sua cultura, fortalecer suas práticas e tradições, ter uma formação pautada em um ensino diferenciado, aqui compreendido como a oportunidade de cultivar e manter as experiências e os modos de vida \textbf{indígenas} nos processos educativos formais.
No entanto, devemos oportunizar, também, a essas/es estudantes o conhecimento sobre outras culturas e outras práticas que permita às/aos \textbf{indígenas} a capacidade de navegar no mundo não \textbf{indígena} e exercer sua cidadania na luta por direitos e no exercício de seus deveres.
Além disso, é preciso investir na formação das/os cidadãs/ãos não \textbf{indígenas}, que respeitem e valorizem os \textbf{indígenas}, que reconheçam suas contribuições na construção da cultura e da \textbf{identidade} brasileira, que sejam parceiros na luta e na garantia dos direitos conquistados por esses povos.
A escola não pode se aceitar como instituição que valida o \textbf{racismo} ou preconceito contra essas/es alunas/os, tampouco pode limitar as capacidades de aprendizagem e, assim, prejudicar seu desenvolvimento intelectual e sociocultural.
Dessa forma, para atender às necessidades educacionais e socioculturais das/as/os alunas/os \textbf{indígenas}, devemos continuar assegurando o direito a uma educação diferenciada e de qualidade, que possa proporcionar o respeito e a preservação de sua cultura, ao mesmo tempo que a/o prepare para o mercado de trabalho, para os estudos posteriores e para o exercício da cidadania, nas sociedades \textbf{indígenas} e não \textbf{indígenas}.

% matched lemmas: carregar, estado, fronteira, hierarquizar, homem, identidade, indígena, mulher, população, racial, racismo, territorialidade, território, índio
\end{textsample}
